\chapter{2021年全国乙卷（理）}

\section{单选题}

\begin{problem}
设 \(2(z+\bar z)+3(z-\bar z)=4+6\i\)，则 \(z=\)（\quad）

\choices
    {\(1-2\i\)}
    {\(1+2\i\)}
    {\(1+\i\)}
    {\(1-\i\)}

\end{problem}

\begin{answer}
C
\end{answer}

\begin{solution}
设 \(z=a+b\i\)，则 \(\bar z=a-b\i\)，
\(2(z+\bar z)+3(z-\bar z)=4a+6b\i=4+6\i\).
所以 \(a=1\)，\(b=1\)，所以 \(z=1+\i\).
\end{solution}
\begin{problem}
已知集合 \(S=\{s\mid s=2n+1,n\in\Z\}\)，\(T=\{t\mid t=4n+1,n\in\Z\}\)，则 \(S\cap T=\)（\quad）

\choices
    {\(\varnothing\)}
    {\(S\)}
    {\(T\)}
    {\(\Z\)}

\end{problem}

\begin{answer}
C
\end{answer}

\begin{solution}
\(s=2n+1\)，\(n\in \Z\).

当 \(n=2k\)，\(k\in \Z\) 时，\(S=\{s\mid s=4k+1,k\in \Z\}\)；当 \(n=2k+1\)，\(k\in \Z\) 时，
\(S=\{s\mid s=4k+3,k\in \Z\}\).
所以 \(T\subseteq S\)，\(S\cap T=T\).故选 \(C\).
\end{solution}
\begin{problem}
已知命题 \(p:\exists x\in \R\)，\(\sin x<1\)；命题 \(q:\forall x\in \R\)，\(\e^{|x|}\ge 1\)，则下列命题中为真命题的是（\quad）

\choices
    {\(p\wedge q\)}
    {\(\neg p\wedge q\)}
    {\(p\wedge \neg q\)}
    {\(\neg\left( p\vee q \right)\)}

\end{problem}

\begin{answer}
A
\end{answer}

\begin{solution}
根据正弦函数的值域 \(\sin x\in [-1,1]\)，故 \(\exists x\in \R\)，\(\sin x<1\)，\(p\) 为真命题.

而函数 \(y=\e^{|x|}\) 为偶函数，且 \(x\ge 0\) 时，\(y=\e^{|x|}\ge 1\)，故 \(\forall x\in \R\)，\(y=\e^{|x|}\ge 1\) 恒成立，则 \(q\) 也为真命题.

所以 \(p\wedge q\) 为真，选 \(A\).
\end{solution}
\begin{problem}
设函数 \(f(x)=\frac{1-x}{1+x}\)，则下列函数中为奇函数的是（\quad）

\choices
    {\(f(x-1)-1\)}
    {\(f(x-1)+1\)}
    {\(f(x+1)-1\)}
    {\(f(x+1)+1\)}

\end{problem}

\begin{answer}
B
\end{answer}

\begin{solution}
\(f(x)=\frac{1-x}{1+x}=-1+\frac{2}{1+x}\).
\(f(x)\) 向右平移一个单位，向上平移一个单位得到
\(g(x)=\frac{2}{x}\)，
为奇函数.
\end{solution}
\begin{problem}
在正方体 \(ABCD-A_1B_1C_1D_1\) 中，\(P\) 为 \(B_1D_1\) 的中点，则直线 \(PB\) 与 \(AD_1\) 所成的角为（\quad）

\bitmapfigure[max width=6.1cm,max totalheight=4.8cm]{img/2021/national_paper_B_science/q5_fig1.png}

\choices
    {\(\frac{\pi}{2}\)}
    {\(\frac{\pi}{3}\)}
    {\(\frac{\pi}{4}\)}
    {\(\frac{\pi}{6}\)}

\end{problem}

\begin{answer}
D
\end{answer}

\begin{solution}
如图，\(\angle PBC_1\) 为直线 \(PB\) 与 \(AD_1\) 所成角的平面角.

在正方体中，\(A_1B\)，\(BC_1\)，\(A_1C_1\) 都是面正方形的对角线，且棱长相等，所以
\(A_1B=BC_1=A_1C_1\)，
故 \(\triangle A_1BC_1\) 为正三角形. 又 \(P\) 为 \(A_1C_1\) 中点，所以 \(BP\) 平分 \(\angle A_1BC_1\)，从而
\(\angle PBC_1=\frac{\pi}{6}\).
\end{solution}
\begin{problem}
将 5 名北京冬奥会志愿者分配到花样滑冰、短道速滑、冰球和冰壶 4 个项目进行培训，每名志愿者只分配到 1 个项目，每个项目至少分配 1 名志愿者，则不同的分配方案共有（\quad）

\choices
    {60 种}
    {120 种}
    {240 种}
    {480 种}

\end{problem}

\begin{answer}
C
\end{answer}

\begin{solution}
所求分配方案数为
\(C_5^2A_4^4=240\).
\end{solution}
\begin{problem}
把函数 \(y=f(x)\) 图像上所有点的横坐标缩短到原来的 \(\frac{1}{2}\) 倍，纵坐标不变，再把所得曲线向右平移 \(\frac{\pi}{3}\) 个单位长度，得到函数 \(y=\sin\left( x-\frac{\pi}{4} \right)\) 的图像，则 \(f(x)=\)（\quad）

\choices
    {\(\sin\left( -\frac{x}{2}-\frac{7\pi}{12} \right)\)}
    {\(\sin\left( \frac{x}{2}+\frac{\pi}{12} \right)\)}
    {\(\sin\left( 2x-\frac{7\pi}{12} \right)\)}
    {\(\sin\left( 2x+\frac{\pi}{12} \right)\)}

\end{problem}

\begin{answer}
B
\end{answer}

\begin{solution}
逆向：\(y=\sin\left( x-\frac{\pi}{4} \right)\xrightarrow{\text{左移 }\frac{\pi}{3}}y=\sin\left( x+\frac{\pi}{12} \right)\xrightarrow{\text{横坐标变为原来的 2 倍}}y=\sin\left( \frac{x}{2}+\frac{\pi}{12} \right)\).

故选 \(B\).
\end{solution}
\begin{problem}
在区间 \((0,1)\) 与 \((1,2)\) 中各随机取 1 个数，则两数之和大于 \(\frac{7}{4}\) 的概率为（\quad）

\choices
    {\(\frac{7}{9}\)}
    {\(\frac{23}{32}\)}
    {\(\frac{9}{32}\)}
    {\(\frac{2}{9}\)}

\end{problem}

\begin{answer}
B
\end{answer}

\begin{solution}
由题意记 \(x\in (0,1)\)，\(y\in (1,2)\)，题目即求 \(x+y>\frac{7}{4}\) 的概率，绘图如下所示.

\(P=\frac{S_{\text{阴}}}{S_{\text{正方形}}} =\frac{1\times 1-\frac{1}{2}\times \frac{3}{4}\times \frac{3}{4}}{1\times 1} =\frac{23}{32}\).
\end{solution}
\begin{problem}
魏晋时期刘徽撰写的《海岛算经》是关于测量的数学著作.其中第一题是测量海岛的高.

如图，点 \(E,H,G\) 在水平线 \(AC\) 上，\(DE\) 和 \(FG\) 是两个垂直于水平面且等高的测量标杆的高度，称为“表高”，\(EG\) 称为“表距”，\(GC\) 和 \(EH\) 都称为“表目距”.\(GC\) 与 \(EH\) 的差称为“表目距的差”，则海岛的高 \(AB=\)（\quad）

\bitmapfigure[max width=5.2cm,max totalheight=4.8cm]{img/2021/national_paper_B_science/q9_fig1.png}

\choices
    {\(\frac{\text{表高}\times \text{表距}}{\text{表目距的差}}+\text{表高}\)}
    {\(\frac{\text{表高}\times \text{表距}}{\text{表目距的差}}-\text{表高}\)}
    {\(\frac{\text{表高}\times \text{表距}}{\text{表目距的差}}+\text{表距}\)}
    {\(\frac{\text{表高}\times \text{表距}}{\text{表目距的差}}-\text{表距}\)}

\end{problem}

\begin{answer}
A
\end{answer}

\begin{solution}
连接 \(DF\) 交 \(AB\) 于 \(M\)，则 \(AB=AM+BM\).

记 \(\angle BDM=\alpha\)，\(\angle BFM=\beta\)，则
\(\frac{MB}{\tan\beta}-\frac{MB}{\tan\alpha}=MF-MD=DF\).

而
\(\tan\beta=\frac{FG}{GC}\)，\(\tan\alpha=\frac{ED}{EH}\).
所以
\(\frac{MB}{\tan\beta}-\frac{MB}{\tan\alpha}=MB\left( \frac{1}{\tan\beta}-\frac{1}{\tan\alpha} \right)=MB\left( \frac{GC}{FG}-\frac{EH}{ED} \right)=MB\cdot \frac{GC-EH}{ED}\).

故
\(MB=\frac{ED\cdot DF}{GC-EH}=\frac{\text{表高}\times \text{表距}}{\text{表目距的差}}\)，
所以高
\(AB=\frac{\text{表高}\times \text{表距}}{\text{表目距的差}}+\text{表高}\).
\end{solution}
\begin{problem}
设 \(a\ne 0\)，若 \(x=a\) 为函数 \(f(x)=a(x-a)^2(x-b)\) 的极大值点，则

\choices
    {\(a<b\)}
    {\(a>b\)}
    {\(ab<a^2\)}
    {\(ab>a^2\)}

\end{problem}

\begin{answer}
D
\end{answer}

\begin{solution}
若 \(a>0\)，其图像如图（1），此时，\(0<a<b\).

\bitmapfigure[max width=4.8cm,max totalheight=4.8cm]{img/2021/national_paper_B_science/q10_solution_fig1.png}

若 \(a<0\)，时图像如图（2），此时，\(b<a<0\).

综上，\(ab<a^2\).
\end{solution}
\begin{problem}
设 \(B\) 是椭圆 \(C:\frac{x^2}{a^2}+\frac{y^2}{b^2}=1\)（\(a>b>0\)）的上顶点，若 \(C\) 上的任意一点 \(P\) 都满足 \(|PB|\le 2b\)，则 \(C\) 的离心率的取值范围是（\quad）

\choices
    {\(\left[ \frac{\sqrt2}{2},1 \right)\)}
    {\(\left[ \frac{1}{2},1 \right)\)}
    {\(\left( 0,\frac{\sqrt2}{2} \right]\)}
    {\(\left( 0,\frac{1}{2} \right]\)}

\end{problem}

\begin{answer}
C
\end{answer}

\begin{solution}
由题意，点 \(B(0,b)\)，设 \(P(x_0,y_0)\)，则
\(\frac{x_0^2}{a^2}+\frac{y_0^2}{b^2}=1\Rightarrow x_0^2=a^2\left( 1-\frac{y_0^2}{b^2} \right)\)，
故
\examdisplayaligned{
|PB|^2&=x_0^2+\left( y_0-b \right)^2\\
&=a^2\left( 1-\frac{y_0^2}{b^2} \right)+y_0^2-2by_0+b^2\\
&=-\frac{c^2}{b^2}y_0^2-2by_0+a^2+b^2,
}
其中 \(y_0\in [-b,b]\).

由题意，当 \(y_0=-b\) 时，\(|PB|^2\) 最大，则
\(-\frac{b^3}{c^2}\le -b\)，\(b^2\ge c^2\)，\(a^2-c^2\ge c^2\).
所以
\(e=\frac{c}{a}\le \frac{\sqrt2}{2}\)，
故
\(e\in \left( 0,\frac{\sqrt2}{2} \right]\).
\end{solution}
\begin{problem}
设 \(a=2\ln 1.01\)，\(b=\ln 1.02\)，\(c=\sqrt{1.04}-1\)，则（\quad）

\choices
    {\(a<b<c\)}
    {\(b<c<a\)}
    {\(b<a<c\)}
    {\(c<a<b\)}

\end{problem}

\begin{answer}
B
\end{answer}

\begin{solution}
设
\(f(x)=\ln(1+x)-\sqrt{1+2x}+1\)，
则 \(b-c=f(0.02)\)，易得
\(f'(x)=\frac{1}{1+x}-\frac{2}{2\sqrt{1+2x}}=\frac{\sqrt{1+2x}-(1+x)}{(1+x)\sqrt{1+2x}}\).

当 \(x\ge 0\) 时，
\(1+x=\sqrt{(1+x)^2}\ge \sqrt{1+2x}\)，
故 \(f'(x)\le 0\).

所以 \(f(x)\) 在 \([0,+\infty)\) 上单调递减，所以 \(f(0.02)<f(0)=0\)，故 \(b<c\).

再设
\(g(x)=2\ln(1+x)-\sqrt{1+4x}+1\)，
则 \(a-c=g(0.01)\)，易得
\(g'(x)=\frac{2}{1+x}-\frac{4}{2\sqrt{1+4x}}=2\cdot \frac{\sqrt{1+4x}-(1+x)}{(1+x)\sqrt{1+4x}}\).

当 \(0\le x<2\) 时，
\(\sqrt{1+4x}\ge \sqrt{1+2x+x^2}=1+x\)，
所以 \(g'(x)\) 在 \([0,2)\) 上 \(\ge 0\).

故 \(g(x)\) 在 \([0,2)\) 上单调递增，所以 \(g(0.01)>g(0)=0\)，故 \(a>c\).

综上，\(a>c>b\).
\end{solution}
\section{填空题}

\begin{problem}
已知双曲线 \(C:\frac{x^2}{m}-y^2=1\)（\(m>0\)）的一条渐近线为 \(\sqrt3 x+my=0\)，则 \(C\) 的焦距为\fillinblank{}.

\end{problem}

\begin{answer}
4
\end{answer}

\begin{solution}
将双曲线 \(C\) 写成
\(\frac{x^2}{a^2}-\frac{y^2}{b^2}=1\)，
则 \(a^2=m\)，\(b^2=1\)，所以双曲线渐近线方程为
\(y=\pm \frac{b}{a}x\)，
由题意，一条渐近线方程为
\(y=-\frac{\sqrt3}{m}x\)，
于是
\(\frac{b}{a}=\frac{1}{\sqrt m}=\frac{\sqrt3}{m}\).
解得 \(m=0\)（舍去）或 \(m=3\)，故焦距为
\(2c=2\sqrt{a^2+b^2}=2\sqrt{3+1}=4\).
\end{solution}
\begin{problem}
已知向量 \(\bs a=(1,3)\)，\(\bs b=(3,4)\)，若 \((\bs a-\lambda \bs b)\perp \bs b\)，则 \(\lambda=\fillinblank{}\).

\end{problem}

\begin{answer}
\(\frac{3}{5}\)
\end{answer}

\begin{solution}
由题意得
\((\bs a-\lambda \bs b)\cdot \bs b=0\)，
即
\(15-25\lambda=0\)，
解得
\(\lambda=\frac{3}{5}\).
\end{solution}
\begin{problem}
记 \(\triangle ABC\) 的内角 \(A\)，\(B\)，\(C\) 的对边分别为 \(a\)，\(b\)，\(c\)，面积为 \(\sqrt3\)，\(B=60^\circ\)，\(a^2+c^2=3ac\)，则 \(b=\fillinblank{}\).

\end{problem}

\begin{answer}
\(2\sqrt2\)
\end{answer}

\begin{solution}
\(S_{\triangle ABC}=\frac{1}{2}ac\sin B=\frac{\sqrt3}{4}ac=\sqrt3\)，
所以 \(ac=4\).

由余弦定理，
\(b^2=a^2+c^2-ac=3ac-ac=2ac=8\)，
所以
\(b=2\sqrt2\).
\end{solution}
\begin{problem}
以图\(\circled{1}\)为正视图，在图\(\circled{2}\)\(\circled{3}\)\(\circled{4}\)\(\circled{5}\)中选两个分别作为侧视图和俯视图，组成某个三棱锥的三视图，则所选侧视图和俯视图的编号依次为\fillinblank{}（写出符合要求的一组答案即可）.

\bitmapfigure[max width=8.2cm,max totalheight=4.8cm]{img/2021/national_paper_B_science/q16_fig1.png}

\end{problem}

\begin{answer}
\(\circled{2}\circled{5}\) 或 \(\circled{3}\circled{4}\)
\end{answer}

\begin{solution}
由高度可知，侧视图只能为 \(\circled{2}\) 或 \(\circled{3}\).

侧视图为 \(\circled{2}\)，如图（1），平面 \(PAC\perp\) 平面 \(ABC\)，\(PA=PC=\sqrt2\)，\(BA=BC=\sqrt5\)，\(AC=2\)，俯视图为 \(\circled{5}\).

俯视图为 \(\circled{3}\)，如图（2），\(PA\perp\) 平面 \(ABC\)，\(PA=1\)，\(AC=AB=\sqrt5\)，\(BC=2\)，俯视图为 \(\circled{4}\).
\end{solution}
\section{解答题}

\begin{problem}
某厂研制了一种生产高精产品的设备，为检验新设备生产产品的某项指标有无提高，用一台旧设备和一台新设备各生产了 10 件产品，得到产品该项指标数据如下：

旧设备为：\(9.8,10.3,10.0,10.2,9.9,9.8,10.0,10.1,10.2,9.7\)

新设备为：\(10.1,10.4,10.1,10.0,10.1,10.3,10.6,10.5,10.4,10.5\)

旧设备和新设备生产产品的该项指标的样本平均数分别记为 \(x\) 和 \(y\)，样本方差分别记为 \(s_1^2\) 和 \(s_2^2\).

\begin{enumerate}
\item 求 \(x\)，\(y\)，\(s_1^2\)，\(s_2^2\)；

\item 判断新设备生产产品的该项指标的均值较旧设备是否有显著提高（如果 \(\dfrac{y-x}{\sqrt{\left( s_1^2+s_2^2 \right)/10}}\ge 2\)，则认为新设备生产产品的该项指标的均值较旧设备有显著提高，否则不认为有显著提高）.
\end{enumerate}
\end{problem}

\begin{solution}
（1）各项所求值如下所示.

\(x=\frac{1}{10}\left( 9.8+10.3+10.0+10.2+9.9+9.8+10.0+10.1+10.2+9.7 \right)=10.0\)，
\(y=\frac{1}{10}\left( 10.1+10.4+10.1+10.0+10.1+10.3+10.6+10.5+10.4+10.5 \right)=10.3\)，
\examdisplayaligned{
s_1^2=&\frac{1}{10}\left[ \left( 9.7-10.0 \right)^2+2\left( 9.8-10.0 \right)^2+\left( 9.9-10.0 \right)^2 \right.\\
&\left.+2\left( 10.0-10.0 \right)^2+\left( 10.1-10.0 \right)^2+2\left( 10.2-10.0 \right)^2+\left( 10.3-10.0 \right)^2 \right]\\
=&0.036,
}
\examdisplayaligned{
s_2^2=&\frac{1}{10}\left[ \left( 10.0-10.3 \right)^2+3\left( 10.1-10.3 \right)^2+\left( 10.3-10.3 \right)^2 \right.\\
&\left.+2\left( 10.4-10.3 \right)^2+2\left( 10.5-10.3 \right)^2+\left( 10.6-10.3 \right)^2 \right]\\
=&0.04.
}

（2）由（1）中数据得
\(y-x=0.3\)，\(2\sqrt{\frac{s_1^2+s_2^2}{10}}\approx 0.34\).
由 \(0.3<0.34\) 可得
\(y-x<2\sqrt{\frac{s_1^2+s_2^2}{10}}\).
所以不认为新设备生产产品的该项指标的均值较旧设备有显著提高.
\end{solution}
\begin{problem}
如图，四棱锥 \(P-ABCD\) 的底面是矩形，\(PD\perp\) 底面 \(ABCD\)，\(PD=DC=1\)，\(M\) 为 \(BC\) 的中点，且 \(PB\perp AM\).

\begin{enumerate}
\item 求 \(BC\)；

\item 求二面角 \(A-PM-B\) 的正弦值.

\bitmapfigure[max width=4.2cm,max totalheight=4.8cm]{img/2021/national_paper_B_science/q18_fig1.png}
\end{enumerate}
\end{problem}

\begin{solution}
（1）因为 \(PD\perp\) 平面 \(ABCD\)，且矩形 \(ABCD\) 中，\(AD\perp DC\).所以以 \(DA\)，\(DC\)，\(DP\) 分别为 \(x\)，\(y\)，\(z\) 轴正方向，\(D\) 为原点建立空间直角坐标系 \(D-xyz\).设 \(BC=t\)，则
\(A(t,0,0)\)，\(B(t,1,0)\)，\(M\left( \frac{t}{2},1,0 \right)\)，\(P(0,0,1)\).
所以
\(\overrightarrow{PB}=(t,1,-1)\)，\(\overrightarrow{AM}=\left( -\frac{t}{2},1,0 \right)\).
因为 \(PB\perp AM\)，所以
\(\overrightarrow{PB}\cdot \overrightarrow{AM}=-\frac{t^2}{2}+1=0\)，
所以
\(t=\sqrt2\)，
故
\(BC=\sqrt2\).

（2）设平面 \(APM\) 的一个法向量为 \(\bs m=(x,y,z)\)，由于
\(\overrightarrow{AP}=(-\sqrt2,0,1)\)，
则
\examdisplaycases{}{
\bs m\cdot \overrightarrow{AP}=-\sqrt2 x+z=0,\\
\bs m\cdot \overrightarrow{AM}=-\frac{\sqrt2}{2}x+y=0.
}
令 \(x=\sqrt2\)，得 \(\bs m=(\sqrt2,1,2)\).

设平面 \(PMB\) 的一个法向量为 \(\bs n=(x',y',z')\)，则
\examdisplaycases{}{
\bs n\cdot \overrightarrow{CB}=\sqrt2 x'=0,\\
\bs n\cdot \overrightarrow{PB}=\sqrt2 x'+y'-z'=0.
}
令 \(y'=1\)，得 \(\bs n=(0,1,1)\).所以
\(\cos \angle\left( \bs m,\bs n \right)=\frac{\bs m\cdot \bs n}{|\bs m||\bs n|}=\frac{3}{\sqrt{14}}\)，
所以二面角 \(A-PM-B\) 的正弦值为
\(\sqrt{1-\frac{9}{14}}=\frac{\sqrt{70}}{14}\).
\end{solution}
\begin{problem}
记 \(S_n\) 为数列 \(\{a_n\}\) 的前 \(n\) 项和，\(b_n\) 为数列 \(\{S_n\}\) 的前 \(n\) 项积，已知 \( \frac{2}{S_n}+\frac{1}{b_n}=2 \).

\begin{enumerate}
\item 证明：数列 \(\{b_n\}\) 是等差数列；

\item 求 \(\{a_n\}\) 的通项公式.
\end{enumerate}
\end{problem}

\begin{solution}
（1）由已知
\(\frac{2}{S_n}+\frac{1}{b_n}=2\)，
则
\(S_n=\frac{b_n}{b_n-1}\)（\(n\ge 2\)），
又因为
\(b_n=S_nb_{n-1}\)，
所以
\(\frac{b_n}{b_{n-1}}=\frac{b_n}{b_n-1}\)，
即
\(b_n-b_{n-1}=\frac{1}{2}\)（\(n\ge 2\)）.
又 \(b_1=S_1\)，由
\(\frac{2}{S_1}+\frac{1}{b_1}=2\)
得
\(b_1=\frac{3}{2}\).
故 \(\{b_n\}\) 是以 \(\frac{3}{2}\) 为首项，\(\frac{1}{2}\) 为公差的等差数列.

（2）由（1）知
\(b_n=\frac{3}{2}+\frac{1}{2}(n-1)=\frac{n+2}{2}\)，
则
\(\frac{2}{S_n}+\frac{2}{n+2}=2\Rightarrow S_n=\frac{n+2}{n+1}\).
\(n=1\) 时，
\(a_1=S_1=\frac{3}{2}\).
\(n\ge 2\) 时，
\(a_n=S_n-S_{n-1}=\frac{n+2}{n+1}-\frac{n+1}{n}=-\frac{1}{n(n+1)}\).
故 \(a_n=\begin{cases}\frac{3}{2},&n=1,\\-\frac{1}{n(n+1)},&n\ge 2.\end{cases}\).
\end{solution}
\begin{problem}
设函数 \(f(x)=\ln(a-x)\)，已知 \(x=0\) 是函数 \(y=xf(x)\) 的极值点.

\begin{enumerate}
\item 求 \(a\)；

\item 设函数 \(g(x)=\dfrac{x+f(x)}{xf(x)}\)，证明：\(g(x)<1\).
\end{enumerate}
\end{problem}

\begin{solution}
（1）令
\(h(x)=xf(x)=x\ln(a-x)\)，
则
\(h'(x)=\ln(a-x)-\frac{x}{a-x}\).
因为 \(x=0\) 是函数 \(y=xf(x)\) 的极值点，所以
\(h'(0)=0\).
解得
\(a=1\).

（2）由（1）可知
\(f(x)=\ln(1-x)\)，
所以
\(g(x)=\frac{x+f(x)}{xf(x)}=\frac{1}{f(x)}+\frac{1}{x}\).

要证 \(g(x)<1\)，即证
\(\frac{1}{\ln(1-x)}+\frac{1}{x}<1\)（\(x<1,\ x\ne 0\)），
即证
\(\frac{x+(1-x)\ln(1-x)}{x\ln(1-x)}<0\).

当 \(x<0\) 时，\(x\ln(1-x)<0\)；当 \(0<x<1\) 时，\(x\ln(1-x)<0\).

所以只需证明
\(x+(1-x)\ln(1-x)>0\).

令
\(H(x)=x+(1-x)\ln(1-x)\)，
\(H(0)=0+(1-0)\ln(1-0)=0\).
则
\(H'(x)=1-\ln(1-x)-\frac{1}{1-x}(1-x)=-\ln(1-x)\).

（1）当 \(x<0\) 时，易得 \(H'(x)<0\)，则 \(H(x)\) 在 \((-\infty,0)\) 上单调递减.

因为 \(H(0)=0\)，所以 \(H(x)>H(0)=0\)，得证.

（2）当 \(0<x<1\) 时，易得 \(H'(x)>0\)，则 \(H(x)\) 在 \((0,1)\) 上单调递增.

因为 \(H(0)=0\)，所以 \(H(x)>H(0)=0\)，得证.

综上证得 \(g(x)<1\).
\end{solution}
\begin{problem}
已知抛物线 \(C:x^2=2py\)（\(p>0\)）的焦点为 \(F\)，且 \(F\) 与圆 \(M:x^2+(y+4)^2=1\) 上点的距离的最小值为 4.

\begin{enumerate}
\item 求 \(p\)；

\item 若点 \(P\) 在 \(M\) 上，\(PA\)，\(PB\) 是 \(C\) 的两条切线，\(A\)，\(B\) 是切点，求 \(\triangle PAB\) 面积的最大值.
\end{enumerate}
\end{problem}

\begin{solution}
（1）焦点
\(F\left( 0,\frac{p}{2} \right)\)
到圆
\(x^2+(y+4)^2=1\)
的最短距离为
\(\frac{p}{2}+3=4\)，
所以
\(p=2\).

（2）抛物线
\(y=\frac{1}{4}x^2\)，
设 \(A(x_1,y_1)\)，\(B(x_2,y_2)\)，\(P(x_0,y_0)\)，得
\(l_{PA}:y=\frac{1}{2}x_1(x-x_1)+y_1=\frac{1}{2}x_1x-\frac{1}{4}x_1^2=\frac{1}{2}x_1x-y_1\)，
\(l_{PB}:y=\frac{1}{2}x_2x-y_2\)，
且
\(x_0^2=-y_0^2-8y_0-15\).

\(l_{PA}\)，\(l_{PB}\) 都过点 \(P(x_0,y_0)\)，则
\examdisplaycases{}{
y_0=\frac{1}{2}x_1x_0-y_1,\\
y_0=\frac{1}{2}x_2x_0-y_2.
}
故
\(l_{AB}:y_0=\frac{1}{2}x_0x-y\)，
即
\(y=\frac{1}{2}x_0x-y_0\).

联立
\examdisplaycases{}{
y=\frac{1}{2}x_0x-y_0,\\
x^2=4y,
}
得
\(x^2-2x_0x+4y_0=0\)，\(\Delta=4x_0^2-16y_0\).

所以
\(|AB|=\sqrt{1+\frac{x_0^2}{4}}\cdot \sqrt{4x_0^2-16y_0}=\sqrt{4+x_0^2}\cdot \sqrt{x_0^2-4y_0}\)，
\(d_{P\to AB}=\frac{|x_0^2-4y_0|}{\sqrt{x_0^2+4}}\)，
所以
\examdisplayaligned{
S_{\triangle PAB}
&=\frac{1}{2}|AB|\cdot d_{P\to AB}\\
&=\frac{1}{2}|x_0^2-4y_0|\cdot \sqrt{x_0^2-4y_0}\\
&=\frac{1}{2}\left( x_0^2-4y_0 \right)^{\frac{3}{2}}\\
&=\frac{1}{2}\left( -y_0^2-12y_0-15 \right)^{\frac{3}{2}}.
}

而 \(y_0\in [-5,-3]\)，故当 \(y_0=-5\) 时，\(S_{\triangle PAB}\) 达到最大，最大值为
\(20\sqrt5\).
\end{solution}
\begin{problem}
在直角坐标系 \(xOy\) 中，圆 \(C\) 的圆心为 \(C(2,1)\)，半径为 1.

\begin{enumerate}
\item 写出圆 \(C\) 的一个参数方程；

\item 过点 \(F(4,1)\) 作圆 \(C\) 的两条切线.以坐标原点为极点，\(x\) 轴正半轴为极轴建立坐标系，求这两条切线的极坐标方程.
\end{enumerate}
\end{problem}

\begin{solution}
（1）圆 \(C\) 的参数方程为 \(x=2+\cos\theta\)，\(y=1+\sin\theta\)（\(\theta\) 为参数）.

（2）圆 \(C\) 的方程为
\((x-2)^2+(y-1)^2=1\).

\circled{1} 当直线斜率不存在时，直线方程为 \(x=4\)，此时圆心到直线距离为 \(2>r\)，舍去；

\circled{2} 当直线斜率存在时，设直线方程为
\(y-1=k(x-4)\)，
化简为
\(kx-y-4k+1=0\)，
此时圆心 \(C(2,1)\) 到直线的距离为
\(d=\frac{|2k-1-4k+1|}{\sqrt{k^2+1}}=r=1\)，
化简得
\(2|k|=\sqrt{k^2+1}\).
两边平方有
\(4k^2=k^2+1\)，
所以
\(k=\pm \frac{\sqrt3}{3}\).

代入直线方程并化简得
\(x-\sqrt3 y+\sqrt3-4=0\)
或
\(x+\sqrt3 y-\sqrt3-4=0\).
化为极坐标方程为
\(\rho \cos\theta-\sqrt3 \rho \sin\theta=4-\sqrt3 \iff \rho \sin\left( \theta+\frac{5\pi}{6} \right)=4-\sqrt3\)，
或
\(\rho \cos\theta+\sqrt3 \rho \sin\theta=4+\sqrt3 \iff \rho \sin\left( \theta+\frac{\pi}{6} \right)=4+\sqrt3\).
\end{solution}
\begin{problem}
已知函数 \(f(x)=|x-a|+|x+3|\).

\begin{enumerate}
\item 当 \(a=1\) 时，求不等式 \(f(x)\ge 6\) 的解集；

\item 若 \(f(x)>-a\)，求 \(a\) 的取值范围.
\end{enumerate}
\end{problem}

\begin{solution}
（1）当 \(a=1\) 时，
\(f(x)\ge 6\iff |x-1|+|x+3|\ge 6\).

当 \(x\le -3\) 时，不等式化为
\(1-x-x-3\ge 6\)，
解得
\(x\le -4\).

当 \(-3<x<1\) 时，不等式化为
\(1-x+x+3\ge 6\)，
解得
\(x\in \varnothing\).

当 \(x\ge 1\) 时，不等式化为
\(x-1+x+3\ge 6\)，
解得
\(x\ge 2\).

综上，原不等式的解集为
\((-\infty,-4]\cup [2,+\infty)\).

（2）若 \(f(x)>-a\)，即 \(f(x)_{\min}>-a\).

因为
\(f(x)=|x-a|+|x+3|\ge |(x-a)-(x+3)|=|a+3|\)
（当且仅当 \((x-a)(x+3)\le 0\) 时，等号成立），所以
\(f(x)_{\min}=|a+3|\)，
所以
\(|a+3|>-a\)，
即
\(a+3<a\)
或
\(a+3>-a\)，
解得
\(a\in \left( -\frac{3}{2},+\infty \right)\).
\end{solution}
